\BeleskeID{03-s014}
% element: 03-s014:e01
% element: 03-s014:e02
% element: 03-s014:e03
% element: 03-s014:e04
% element: 03-s014:d01

Kada su selektori povezani sa promenljivama funkcije, podatkovni ulaz indeksiran adresom \(BA\) dobija vrednost \(F(B,A)\) iz odgovarajućeg reda funkcionalne tabele. Redosled 00, 01, 10, 11 daje vrednosti 1, 0, 0, 1. Zato postavljene konstante ne predstavljaju nove promenljive: one zadaju šta multiplekser treba da vrati za svaku adresu.

Uvrštavanjem u izraz za MUX dobija se
\[
 Z=1\cdot\overline B\overline A
   +0\cdot\overline BA
   +0\cdot B\overline A+1\cdot BA
   =\overline B\overline A+BA=F.
\]
Dekoder formira minterme, pa se potrebni izlazi naknadno sabiraju. Multiplekser već sadrži uslovljavanje podatkovnih grana i njihovo ILI objedinjavanje; u ovom primeru dovoljno je povezati selektore i konstantne podatke. Zbog toga realizacija multiplekserom može biti jednostavnija.
